\documentclass[12pt,a4paper]{article}

\usepackage[T2A]{fontenc}
\usepackage[utf8]{inputenc}
\usepackage[russian]{babel}
\usepackage{amsmath,amssymb,amsthm}
\usepackage{geometry}
\usepackage{hyperref}
\usepackage{enumitem}

\geometry{left=25mm,right=20mm,top=20mm,bottom=20mm}

\newcommand{\R}{\mathbb{R}}
\newcommand{\eps}{\varepsilon}
\DeclareMathOperator{\sgn}{sgn}
\newcommand{\source}[1]{\par\noindent\textit{Источник:} #1.}

\theoremstyle{definition}
\newtheorem{definition}{Определение}[section]
\newtheorem{example}{Пример}[section]

\theoremstyle{plain}
\newtheorem{theorem}{Теорема}[section]
\newtheorem{lemma}{Лемма}[section]
\newtheorem{corollary}{Следствие}[section]

\theoremstyle{remark}
\newtheorem{remark}{Замечание}[section]

\begin{document}

% Титульный лист вставьте здесь.

\tableofcontents
\newpage

\section*{Введение}
\addcontentsline{toc}{section}{Введение}

Несобственные интегралы возникают при распространении определённого интеграла на бесконечные промежутки и на функции, неограниченные в окрестности особой точки. Такие интегралы играют важную роль в анализе, теории вероятностей, математической физике, теории специальных функций. Например, интеграл Пуассона
\[
\int_0^\infty e^{-x^2}\,dx=\frac{\sqrt{\pi}}{2}
\]
лежит в основе нормального закона распределения, а интеграл Дирихле
\[
\int_0^\infty \frac{\sin x}{x}\,dx=\frac{\pi}{2}
\]
является классическим примером условно сходящегося интеграла и часто используется при вычислении интегралов через параметр.

\textbf{Цель работы} — систематизация признаков сходимости несобственных интегралов, освоение методов их исследования и вычисления.

\textbf{Задачи работы:}
\begin{enumerate}[label=\arabic*)]
\item дать определения несобственных интегралов I и II рода;
\item сформулировать и доказать основные признаки сходимости;
\item изучить методы вычисления несобственных интегралов, включая сведение к специальным функциям и дифференцирование по параметру;
\item разобрать содержательные примеры, в том числе условно сходящиеся интегралы.
\end{enumerate}

Работа опирается прежде всего на учебники Г.~М.~Фихтенгольца \cite{fikht} и В.~А.~Зорича \cite{zorich}. Ссылки на страницы указываются по мере изложения.

\section{Основные понятия и определения}

\subsection{Несобственные интегралы I рода}
\label{sec:I}

\begin{definition}[Несобственный интеграл I рода]
Пусть функция \(f\) определена на промежутке \([a,+\infty)\) и интегрируема на каждом конечном отрезке \([a,A]\), \(A>a\). Если существует конечный предел
\[
\lim_{A\to+\infty}\int_a^A f(x)\,dx,
\]
то этот предел называется несобственным интегралом I рода от \(f\) по \([a,+\infty)\) и обозначается
\[
\int_a^{+\infty}f(x)\,dx.
\]
В этом случае интеграл называется сходящимся. Если предел не существует или бесконечен, интеграл называется расходящимся.
\end{definition}
\source{Фихтенгольц, т.~2, с.~552--553}

Аналогично определяются интегралы
\[
\int_{-\infty}^b f(x)\,dx=\lim_{A'\to-\infty}\int_{A'}^b f(x)\,dx,
\]
\[
\int_{-\infty}^{+\infty}f(x)\,dx
=
\int_{-\infty}^c f(x)\,dx+\int_c^{+\infty}f(x)\,dx,
\]
если оба интеграла справа сходятся. Выбор точки \(c\) не влияет на сходимость.

\begin{example}
\[
\int_1^{+\infty}\frac{dx}{x^\lambda}
\]
сходится при \(\lambda>1\) и расходится при \(\lambda\le 1\). В самом деле, при \(\lambda\ne1\)
\[
\int_1^A x^{-\lambda}\,dx=\frac{A^{1-\lambda}-1}{1-\lambda},
\]
и предел при \(A\to+\infty\) конечен только при \(\lambda>1\). При \(\lambda=1\)
\[
\int_1^A\frac{dx}{x}=\ln A\to+\infty.
\]
\end{example}
\source{Фихтенгольц, т.~2, с.~553}

Если для \(f\) существует первообразная \(F\) на \([a,+\infty)\), то
\[
\int_a^{+\infty}f(x)\,dx=F(+\infty)-F(a),
\]
где \(F(+\infty)=\lim_{x\to+\infty}F(x)\), при условии существования конечного предела. Эту запись удобно понимать как обобщённую формулу Ньютона—Лейбница.
\source{Фихтенгольц, т.~2, с.~554}

\subsection{Несобственные интегралы II рода}
\label{sec:II}

\begin{definition}[Несобственный интеграл II рода]
Пусть функция \(f\) определена на промежутке \([a,b)\), интегрируема на каждом отрезке \([a,b']\), где \(a\le b'<b\), и неограничена в окрестности точки \(b\) (точка \(b\) — особая). Если существует конечный предел
\[
\lim_{b'\to b-0}\int_a^{b'}f(x)\,dx,
\]
то он называется несобственным интегралом II рода от \(f\) по \([a,b]\) и обозначается
\[
\int_a^b f(x)\,dx.
\]
Если предел конечен, интеграл сходится; иначе расходится.
\end{definition}
\source{Фихтенгольц, т.~2, с.~577--579}

Аналогично, если особая точка находится в левом конце \(a\), то
\[
\int_a^b f(x)\,dx=\lim_{a'\to a+0}\int_{a'}^b f(x)\,dx.
\]
Если особая точка \(c\in(a,b)\), то интеграл определяется как сумма
\[
\int_a^b f(x)\,dx=\int_a^c f(x)\,dx+\int_c^b f(x)\,dx
\]
при условии сходимости обоих интегралов справа.

\begin{example}
\[
\int_a^b\frac{dx}{(x-a)^\lambda}\qquad (b>a)
\]
сходится при \(\lambda<1\) и расходится при \(\lambda\ge1\). Действительно, при \(\lambda\ne1\)
\[
\int_{a+\eta}^b (x-a)^{-\lambda}\,dx
=
\frac{(b-a)^{1-\lambda}-\eta^{1-\lambda}}{1-\lambda},
\]
и предел при \(\eta\to+0\) конечен только при \(\lambda<1\). При \(\lambda=1\) получаем \(\ln(b-a)-\ln\eta\to+\infty\).
\end{example}
\source{Фихтенгольц, т.~2, с.~579--580}

\subsection{Абсолютная и условная сходимость}

\begin{definition}
Несобственный интеграл \(\int_a^\omega f(x)\,dx\) называется абсолютно сходящимся, если сходится интеграл
\[
\int_a^\omega |f(x)|\,dx.
\]
Если \(\int_a^\omega f\) сходится, а \(\int_a^\omega |f|\) расходится, то интеграл называется условно сходящимся.
\end{definition}
\source{Фихтенгольц, т.~2, с.~562--563; 584--586}

\begin{theorem}[Абсолютная сходимость влечёт сходимость]
Если интеграл \(\int_a^\omega |f(x)|\,dx\) сходится, то интеграл \(\int_a^\omega f(x)\,dx\) также сходится, причём
\[
\left|\int_a^\omega f(x)\,dx\right|\le \int_a^\omega |f(x)|\,dx.
\]
\end{theorem}

\begin{proof}
Для интеграла I рода применим критерий Коши к \(\int_a^\infty |f|\). Для любых \(A'>A>A_0\)
\[
\left|\int_A^{A'} f(x)\,dx\right|
\le
\int_A^{A'} |f(x)|\,dx<\eps.
\]
Значит, для \(\int_a^\infty f\) выполнен критерий Коши, и он сходится. Для II рода рассуждение аналогично с пределами \(b'\to b-0\).
\end{proof}
\source{Фихтенгольц, т.~2, с.~562}

\begin{example}
Интеграл
\[
\int_1^{+\infty}\frac{\sin x}{x}\,dx
\]
сходится по признаку Дирихле, но не абсолютно. Неабсолютность следует из расходимости \(\int_1^\infty |\sin x|/x\,dx\): на отрезках \([k\pi,(k+1)\pi]\)
\[
\int_{k\pi}^{(k+1)\pi}\frac{|\sin x|}{x}\,dx
\ge
\frac{1}{(k+1)\pi}\int_0^\pi \sin u\,du
=
\frac{2}{(k+1)\pi},
\]
а гармонический ряд расходится.
\end{example}
\source{Фихтенгольц, т.~2, с.~565}

\section{Признаки сходимости несобственных интегралов}

\subsection{Интегральный признак Коши для положительных функций}

\begin{theorem}[Признак Коши]
Пусть \(f(x)\ge0\) на \([a,+\infty)\). Тогда интеграл \(\int_a^{+\infty}f(x)\,dx\) сходится в том и только том случае, когда функция
\[
F(A)=\int_a^A f(x)\,dx
\]
ограничена сверху на \([a,+\infty)\).
\end{theorem}

\begin{proof}
Так как \(f\ge0\), функция \(F(A)\) монотонно не убывает. По теореме о пределе монотонной функции она имеет конечный предел при \(A\to+\infty\) тогда и только тогда, когда ограничена сверху.
\end{proof}
\source{Фихтенгольц, т.~2, с.~559--560}

\begin{corollary}[Сравнение с \(1/x^\lambda\)]
Пусть \(f(x)\ge0\) и при достаточно больших \(x\)
\[
f(x)=\frac{\varphi(x)}{x^\lambda},\qquad \lambda>0.
\]
\begin{enumerate}
\item Если \(\lambda>1\) и \(0\le\varphi(x)\le C\), то \(\int_a^\infty f(x)\,dx\) сходится.
\item Если \(\lambda\le1\) и \(\varphi(x)\ge c>0\), то \(\int_a^\infty f(x)\,dx\) расходится.
\end{enumerate}
\end{corollary}
\source{Фихтенгольц, т.~2, с.~560}

Для интегралов II рода аналогично:
\[
f(x)=\frac{\psi(x)}{(b-x)^\lambda}.
\]
Если \(\lambda<1\) и \(0\le\psi(x)\le C\), то \(\int_a^b f\) сходится; если \(\lambda\ge1\) и \(\psi(x)\ge c>0\), то расходится.
\source{Фихтенгольц, т.~2, с.~584--585}

\subsection{Признаки сравнения}

\begin{theorem}[Обычное сравнение]
Пусть \(0\le f(x)\le g(x)\) при \(x\ge A\). Тогда:
\begin{enumerate}
\item из сходимости \(\int_a^\infty g(x)\,dx\) следует сходимость \(\int_a^\infty f(x)\,dx\);
\item из расходимости \(\int_a^\infty f(x)\,dx\) следует расходимость \(\int_a^\infty g(x)\,dx\).
\end{enumerate}
\end{theorem}
\source{Фихтенгольц, т.~2, с.~560}

\begin{theorem}[Предельное сравнение]
Пусть \(f(x)\ge0\), \(g(x)>0\) и существует предел
\[
\lim_{x\to+\infty}\frac{f(x)}{g(x)}=K,\qquad 0\le K\le+\infty.
\]
Тогда:
\begin{enumerate}
\item если \(K<+\infty\) и \(\int_a^\infty g\) сходится, то \(\int_a^\infty f\) сходится;
\item если \(K>0\) и \(\int_a^\infty g\) расходится, то \(\int_a^\infty f\) расходится.
\end{enumerate}
В частности, при \(0<K<+\infty\) интегралы сходятся или расходятся одновременно.
\end{theorem}
\source{Фихтенгольц, т.~2, с.~560}

\subsection{Признаки Дирихле и Абеля}

\begin{theorem}[Признак Дирихле]
Пусть функция \(f\) локально интегрируема на \([a,+\infty)\), а функция \(g\) монотонна на \([a,+\infty)\) и
\[
g(x)\to0\quad (x\to+\infty).
\]
Если частичные интегралы
\[
F(A)=\int_a^A f(x)\,dx
\]
ограничены:
\[
|F(A)|\le M,
\]
то интеграл
\[
\int_a^{+\infty}f(x)g(x)\,dx
\]
сходится.
\end{theorem}

\begin{proof}[Доказательство, схема]
По второй теореме о среднем для любых \(A'>A\) найдётся \(\xi\in[A,A']\) такое, что
\[
\int_A^{A'}f(x)g(x)\,dx
=
g(A)\int_A^\xi f(x)\,dx
+
g(A')\int_\xi^{A'}f(x)\,dx.
\]
Оба интеграла справа по модулю не превосходят \(2M\), поэтому
\[
\left|\int_A^{A'}f g\,dx\right|
\le 4M\max\{|g(A)|,|g(A')|\}.
\]
Так как \(g(x)\to0\), правая часть стремится к нулю при \(A,A'\to+\infty\). По критерию Коши интеграл сходится.
\end{proof}
\source{Фихтенгольц, т.~2, с.~563--565}

\begin{theorem}[Признак Абеля]
Пусть интеграл \(\int_a^{+\infty}f(x)\,dx\) сходится, а функция \(g\) монотонна и ограничена на \([a,+\infty)\). Тогда интеграл \(\int_a^{+\infty}f(x)g(x)\,dx\) сходится.
\end{theorem}

\begin{proof}[Доказательство, схема]
Пусть \(L=\lim_{x\to\infty}g(x)\). Тогда
\[
fg=fL+f(g-L).
\]
Интеграл от \(fL\) сходится, а для \(f(g-L)\) выполнены условия Дирихле, поскольку \(g-L\to0\) монотонно. Значит, сходится и \(\int fg\).
\end{proof}
\source{Фихтенгольц, т.~2, с.~564--565}

\subsection{Критерий Коши сходимости несобственного интеграла}

\begin{theorem}[Критерий Коши для I рода]
Интеграл \(\int_a^{+\infty}f(x)\,dx\) сходится тогда и только тогда, когда для любого \(\eps>0\) существует \(A_0>a\) такое, что для всех \(A',A>A_0\)
\[
\left|\int_A^{A'}f(x)\,dx\right|<\eps.
\]
\end{theorem}
\source{Фихтенгольц, т.~2, с.~562}

\begin{theorem}[Критерий Коши для II рода]
Пусть особая точка — \(b\). Интеграл \(\int_a^b f(x)\,dx\) сходится тогда и только тогда, когда для любого \(\eps>0\) существует \(\delta>0\) такое, что для всех \(a_1,a_2\in(b-\delta,b)\), \(a_1<a_2\),
\[
\left|\int_{a_1}^{a_2}f(x)\,dx\right|<\eps.
\]
\end{theorem}
\source{Фихтенгольц, т.~2, с.~584--586}

\subsection{Эквивалентные преобразования подынтегральной функции}

\begin{corollary}
Пусть \(f(x)\ge0\), \(g(x)>0\) и \(f(x)\sim g(x)\) при \(x\to+\infty\), то есть
\[
\lim_{x\to+\infty}\frac{f(x)}{g(x)}=1.
\]
Тогда интегралы
\[
\int_a^\infty f(x)\,dx,\qquad \int_a^\infty g(x)\,dx
\]
сходятся или расходятся одновременно.
\end{corollary}

Аналогичное утверждение верно для II рода при \(x\to b-0\), если \(f,g\) не меняют знак вблизи особой точки.
\source{Фихтенгольц, т.~2, с.~560}

\section{Методы вычисления несобственных интегралов}

\subsection{Интегрирование по частям и замена переменной}

\begin{theorem}[Интегрирование по частям]
Пусть \(u,v\) непрерывно дифференцируемы на \([a,b)\), где \(b\) может быть равно \(+\infty\). Тогда
\[
\int_a^b u\,dv
=
uv\big|_a^b-\int_a^b v\,du,
\]
если из трёх выражений: двух интегралов и двойной подстановки — имеют смысл любые два. Двойная подстановка понимается как
\[
uv\big|_a^b=\lim_{x\to b-0}u(x)v(x)-u(a)v(a).
\]
\end{theorem}
\source{Фихтенгольц, т.~2, с.~602}

\begin{theorem}[Замена переменной]
Пусть \(x=\varphi(t)\) — монотонно возрастающая непрерывно дифференцируемая функция на \([\alpha,\beta)\), \(\varphi(\alpha)=a\), \(\lim_{t\to\beta-0}\varphi(t)=b\). Тогда
\[
\int_a^b f(x)\,dx=\int_\alpha^\beta f(\varphi(t))\varphi'(t)\,dt,
\]
если существует один из этих интегралов; тогда существует и другой.
\end{theorem}
\source{Фихтенгольц, т.~2, с.~604--605}

\subsection{Сведение к собственному интегралу и предельный переход}

Если удаётся найти первообразную \(F\), то
\[
\int_a^b f(x)\,dx=\lim_{x\to b-0}F(x)-F(a).
\]
Для I рода:
\[
\int_a^{+\infty}f(x)\,dx=\lim_{A\to+\infty}F(A)-F(a).
\]
Это основной вычислительный приём, если первообразная выражается через элементарные функции.
\source{Фихтенгольц, т.~2, с.~554; 582}

\subsection{Использование специальных функций}

\begin{definition}[Гамма-функция Эйлера]
Для \(\alpha>0\) полагают
\[
\Gamma(\alpha)=\int_0^{+\infty}x^{\alpha-1}e^{-x}\,dx.
\]
\end{definition}

\begin{definition}[Бета-функция Эйлера]
Для \(\alpha>0,\ \beta>0\) полагают
\[
B(\alpha,\beta)=\int_0^1 x^{\alpha-1}(1-x)^{\beta-1}\,dx.
\]
\end{definition}

Основные свойства:
\[
\Gamma(\alpha+1)=\alpha\Gamma(\alpha),\qquad \Gamma(n+1)=n!,
\]
\[
B(\alpha,\beta)=\frac{\Gamma(\alpha)\Gamma(\beta)}{\Gamma(\alpha+\beta)}.
\]
В частности,
\[
\Gamma\left(\frac12\right)=\sqrt{\pi}.
\]
\source{Зорич, ч.~II, с.~406--416}

\subsection{Дифференцирование по параметру}

\begin{theorem}[Дифференцирование несобственного интеграла по параметру]
Пусть \(f(x,y)\) и \(f_y(x,y)\) непрерывны на \([a,+\infty)\times[c,d]\). Пусть интеграл
\[
\int_a^{+\infty}f(x,y_0)\,dx
\]
сходится хотя бы при одном \(y_0\in[c,d]\), а интеграл
\[
\int_a^{+\infty}f_y(x,y)\,dx
\]
сходится равномерно по \(y\in[c,d]\). Тогда интеграл
\[
F(y)=\int_a^{+\infty}f(x,y)\,dx
\]
сходится равномерно на \([c,d]\), функция \(F\) дифференцируема и
\[
F'(y)=\int_a^{+\infty}f_y(x,y)\,dx.
\]
\end{theorem}
\source{Зорич, ч.~II, с.~396--398}

\section{Примеры исследования и вычисления}

\subsection{Интеграл Пуассона}

Рассмотрим
\[
I=\int_0^{+\infty}e^{-x^2}\,dx.
\]
Этот интеграл сходится, поскольку при \(x\to+\infty\)
\[
e^{-x^2}=o\left(\frac1{x^2}\right),
\]
а \(\int_1^\infty x^{-2}\,dx\) сходится.

Вычислим \(I\). Имеем
\[
I^2=\int_0^{+\infty}\int_0^{+\infty}e^{-(x^2+y^2)}\,dx\,dy.
\]
Переходя к полярным координатам и используя неотрицательность подынтегральной функции, получаем
\[
I^2=\int_0^{\pi/2}d\varphi\int_0^{+\infty}e^{-r^2}r\,dr
=
\frac{\pi}{2}\cdot\frac12
=
\frac{\pi}{4}.
\]
Следовательно,
\[
I=\frac{\sqrt{\pi}}{2}.
\]
Тот же результат получается из формулы
\[
\int_0^{+\infty}e^{-x^2}\,dx
=
\frac12\Gamma\left(\frac12\right)
=
\frac{\sqrt{\pi}}{2}.
\]
\source{Фихтенгольц, т.~2, с.~612--613; Зорич, ч.~II, с.~402--403, 410--411}

\subsection{Интеграл Дирихле}

Рассмотрим
\[
J=\int_0^{+\infty}\frac{\sin x}{x}\,dx.
\]
В нуле подынтегральная функция имеет конечный предел \(1\), поэтому особенность только на бесконечности. По признаку Дирихле интеграл сходится: \(f(x)=\sin x\), \(g(x)=1/x\), функция \(g\) монотонно стремится к нулю, а интегралы от \(\sin x\) ограничены.

Покажем, что сходимость неабсолютная. На отрезках \([k\pi,(k+1)\pi]\)
\[
\int_{k\pi}^{(k+1)\pi}\frac{|\sin x|}{x}\,dx
\ge
\frac{1}{(k+1)\pi}\int_0^\pi\sin u\,du
=
\frac{2}{(k+1)\pi},
\]
и ряд из этих величин расходится. Значит,
\[
\int_0^{+\infty}\frac{|\sin x|}{x}\,dx=+\infty,
\]
то есть интеграл Дирихле сходится условно.

Вычислим его. Рассмотрим
\[
F(y)=\int_0^{+\infty}e^{-xy}\frac{\sin x}{x}\,dx,\qquad y\ge0.
\]
При \(y>0\) дифференцируем по параметру:
\[
F'(y)=-\int_0^{+\infty}e^{-xy}\sin x\,dx
=
-\frac{1}{1+y^2}.
\]
Кроме того, \(F(y)\to0\) при \(y\to+\infty\). Поэтому
\[
F(y)=\frac{\pi}{2}-\arctan y.
\]
Переходя к пределу при \(y\to+0\) и используя равномерную сходимость (признак Абеля—Дирихле), получаем
\[
J=F(0)=\frac{\pi}{2}.
\]
\source{Фихтенгольц, т.~2, с.~565, 613; Зорич, ч.~II, с.~398}

\section*{Заключение}
\addcontentsline{toc}{section}{Заключение}

В работе систематизированы определения несобственных интегралов I и II рода, основные признаки сходимости: критерий Коши, признаки сравнения, признак Коши для положительных функций, признаки Дирихле и Абеля. Рассмотрены методы вычисления: сведение к пределу собственного интеграла, интегрирование по частям, замена переменной, использование гамма- и бета-функций, дифференцирование по параметру. Разобраны интеграл Пуассона и условно сходящийся интеграл Дирихле. Показано, что выбор признака и метода вычисления существенно зависит от знака подынтегральной функции, наличия абсолютной сходимости и поведения функции вблизи особой точки.

\begin{thebibliography}{9}

\bibitem{fikht}
Фихтенгольц~Г.~М.
Курс дифференциального и интегрального исчисления. Т.~2.
--- М.: Наука, 1969.

\bibitem{zorich}
Зорич~В.~А.
Математический анализ. Часть~II.
--- М.: МЦНМО, 2019.

\end{thebibliography}

\end{document}